\section{Binomial shifts of the meander numbers}\label{sec:idle-moves}

Write \(M(z)=\sum_{n\ge1}M_nz^n\), where \(2n\) is the number of
crossings, and put
\[
     F_t(z)=M\!\left(\frac{z}{1+tz}\right),\qquad
     F_{t,n}=[z^n]F_t(z),\qquad C_n=F_{2,n}.
\]
Thus \(F_{0,n}=M_n\), and the one-sign avoidance identity gives
\(F_{1,n}=B_n\).

\begin{theorem}\label{thm:idle-shifts}
All coefficients of \(F_t\) are nonnegative if and only if \(t\le2\).
For each real \(t\le2\) they are eventually positive and
\[
       \lim_{n\to\infty}\frac{F_{t,n+1}}{F_{t,n}}=\mu-t .
\]
In particular the assertion at \(t=2\) holds despite \(C_2=0\).
Writing \(E M_n=M_{n+1}\), one has
\((E-2)^kM_n\ge0\) for every \(k\ge0\) and \(n\ge1\).
\end{theorem}

\begin{proof}
Use the four-letter sweep of a closed meander
\cite{Jensen,AP}. Turn the fixed line vertical and sweep upwards.
The state records the ordered ends on each side of the line and which
two ends belong to the same unfinished arc. A letter \(U\) creates a
fresh arc with one end on either side. A letter \(D\) joins the two
nearest ends on opposite sides. The letters \(R,L\) move the nearest
end from left to right and from right to left, respectively.
An unfinished word may never close a component. The first letter is
\(U\), giving a state \(e\) with one arc. The last letter is \(D\),
and its preceding state must be \(e\). The choices of letters determine
the two noncrossing matchings, since an end always closes against the
nearest available end on its side.

After an odd number of crossings there is at least one end on either
side. Indeed, on each side the number of ends equals the number of
crossings processed minus twice the number of completed matching edges
on that side. Divide the \(2n-2\) interior letters into consecutive
disjoint pairs, beginning immediately after the first \(U\). Write
\(T_{xy}\) for the number of legal two-letter words taking odd-cut
state \(x\) to state \(y\). Each row has sum at most sixteen, so every
entry of every power used below is a finite sum. Bijectivity of the
sweep gives \(M_{m+1}=(T^m)_{ee}\). At each odd cut the distinct pairs
\(LR\) and \(RL\) are legal and return the original state. Delete these
two choices at every state and call the resulting matrix \(Q=T-2I\).
In particular \(d_m=(Q^m)_{ee}\) counts meanders of order \(m+1\)
with neither pair in the designated positions.

The binomial theorem now gives
\[
   M_{m+1}=\sum_{j=0}^m\binom mj 2^{m-j}d_j,\qquad
   (E-2)^kM_n=(T^{n-1}Q^k)_{ee}.
\]
The first identity gives \(d_m=C_{m+1}\) by binomial inversion, and
the second proves the finite-difference assertion. More generally, for
\(s=2-t\) another binomial expansion gives
\begin{equation}\label{eq:idle-weight}
   F_{t,m+1}=\big((Q+sI)^m\big)_{ee}
            =\sum_{j=0}^m\binom mj s^{m-j}d_j .
\end{equation}
This follows also by substituting \(z/(1-sz)\) in \(F_2\).
It proves nonnegativity when \(s\ge0\). The coefficient
\(F_{t,2}=M_2-tM_1=2-t\) gives the converse.

Concatenating two return words at \(e\) gives
\(d_{m+h}\ge d_md_h\). The first four values are
\(d_0=1,\ d_1=0,\ d_2=4,\ d_3=10\), by binomial inversion from
\(M_1=1,\ M_2=2,\ M_3=8,\ M_4=42\). Every integer \(m\ge2\) is
a sum of twos and threes, so \(d_m>0\) in this range. Put
\(c=\sup_{k\ge2} d_k^{1/k}\le16\). For fixed \(k\ge2\) and all
large \(m\), write \(m=ak+b\) with \(2\le b\le k+1\).
Then \(d_m\ge d_k^a d_b\), so \(d_m^{1/m}\to c\) and \(c\ge2\).
For \(s>0\), the \(m\)-th root of the right side of
\eqref{eq:idle-weight} tends to \(s+c\). The upper bound follows
from \(d_j\le c^j\). For the lower bound use \(d_j\ge(c-\epsilon)^j\)
for all sufficiently large \(j\), with \(0<\epsilon<c\). The finitely
many omitted terms are \(O(m^K s^m)\) for fixed \(K\), which has
smaller exponential growth than \((s+c-\epsilon)^m\). Setting \(s=2\)
identifies \(c=\mu-2\). For any \(s>0\), the sequence
\(a_m=s^{-m}F_{t,m+1}\) is the binomial transform of the nonnegative
sequence \(s^{-j}d_j\) with initial term one, and its root tends to
\(1+c/s>1\). Lemma~\ref{lem:binomial-smoothing} gives
\(F_{t,m+2}/F_{t,m+1}\to s+c=\mu-t\).

It remains to treat \(s=0\). Two longer return words are available
at every odd cut even after the two original pairs have been deleted:
\[
               UR\mid DL,\qquad UU\mid RD\mid DL .
\]
Here the bars mark the disjoint pairs. To verify both words, note that
\(URD\) has exactly the effect of \(R\) whenever a left end is present.
The new arc created by \(U\) is separate from every old component.
The next letter moves its left end to the right. The letter \(D\)
joins that end to the original nearest left end, extending the old
arc by the fresh arc and placing its end on the right. The order and
pairing of the other ends are unchanged. Thus \(URDL\) returns the
original state, as in Albert and Paterson's zero-displacement example
\cite{AP}. In \(U(URD)DL\) the inner \(URD\) acts on the fresh
left end of the outer \(U\), and replacing it by \(R\) leaves \(URDL\).
Both constructions are legal even at \(e\). In particular
\(Q^r-I\) is entrywise nonnegative for \(r=2,3\).

Fix \(r\in\{2,3\}\) and use \(j=2,\ldots,r+1\) as representatives
of all residue classes modulo \(r\). Then
\[
 d_{rm+j}
    =\sum_{\ell=0}^m\binom m\ell
         \big(Q^j(Q^r-I)^\ell\big)_{ee}.
\]
The coefficients are nonnegative, their initial term is \(d_j>0\),
and the \(m\)-th root of the left side tends to \(c^r>1\).
Lemma~\ref{lem:binomial-smoothing} therefore gives
\(d_{n+r}/d_n\to c^r\) along every residue class and hence along
all integers. Finally,
\[
    \frac{d_{n+1}}{d_n}
       =\frac{d_{n+3}/d_n}{d_{n+3}/d_{n+1}}\longrightarrow c .
\]
Since \(d_n=C_{n+1}\), this proves the endpoint assertion.
\end{proof}
